\subsubsection{Fluxo de Bens Patrimoniais: Requisições}

\begin{lstlisting}[language=TypeScript, title={Criação e Resposta de Requisição de Edição (C2: lock determinístico)}]
export const criarRequisicaoEdicaoBem = onCall(async (request) => {
  await validarPermissao(request, ["Professor"]);
  const dados = request.data as { idBemPatrimonial: string; novoNome?: string; novoStatus?: string; novoEstadoConservacao?: string; novoIdLocal?: string; motivo: string };

  // C2: Lock determinístico - elimina phantom reads em transações Firestore.
  // Se dois professores enviarem simultaneamente, apenas um cria o lock; o outro
  // recebe ALREADY_EXISTS da transação e falha antes de criar a requisição.
  const lockId = `bem_edicao_${dados.idBemPatrimonial}`;
  const lockRef = admin.firestore().collection("Locks_Requisicao_Patrimonio").doc(lockId);
  const reqRef  = admin.firestore().collection("Requisicao_Edicao_Bem_Patrimonial").doc();

  return admin.firestore().runTransaction(async (tx) => {
    const lockSnap = await tx.get(lockRef);
    if (lockSnap.exists) {
      throw new HttpsError("failed-precondition", "Já existe uma requisição de edição pendente para este bem.");
    }
    const bemSnap = await tx.get(admin.firestore().collection("Bem_Patrimonial").doc(dados.idBemPatrimonial));
    if (!bemSnap.exists) throw new HttpsError("not-found", "Bem não encontrado.");
    // Cria lock e requisição atomicamente
    tx.set(lockRef, { criado_em: admin.firestore.FieldValue.serverTimestamp() });
    tx.set(reqRef, {
      id_bem_patrimonial: dados.idBemPatrimonial,
      novo_nome: dados.novoNome ?? null,
      versao_bem_origem: bemSnap.data()!.versao,
      novo_id_resumo_bem_patrimonial: null,
      novo_status: dados.novoStatus ?? null,
      novo_estado_conservacao: dados.novoEstadoConservacao ?? null,
      novo_id_local: dados.novoIdLocal ?? null,
      motivo: dados.motivo,
      status: "pendente",
      feita_em: admin.firestore.FieldValue.serverTimestamp(),
      id_usuario_solicitante: request.auth!.uid,
    });
    return { idRequisicao: reqRef.id };
  });
});

export const responderRequisicaoEdicaoBem = onCall(async (request) => {
  await validarPermissao(request, ["Chefe_Geral", "Gestor_Bens_Patrimoniais"], true);
  const { idRequisicao, aprovar, justificativa, idResumoAlvo } = request.data;
  return admin.firestore().runTransaction(async (tx) => {
    const reqRef = admin.firestore().collection("Requisicao_Edicao_Bem_Patrimonial").doc(idRequisicao);
    const req = (await tx.get(reqRef)).data();
    if (!req || req.status !== "pendente") throw new HttpsError("failed-precondition", "Requisição indisponível.");
    const bemRef = admin.firestore().collection("Bem_Patrimonial").doc(req.id_bem_patrimonial);
    const bem = (await tx.get(bemRef)).data();
    // N-02: versao_bem_origem garante que aprovação sobre edição concorrente seja rejeitada.
    // Se o bem foi editado após a abertura da requisição, a versão diverge e a aprovação
    // é abortada -- o respondente deve revisar a proposta com os dados atuais.
    if (!bem || bem.versao !== req.versao_bem_origem)
      throw new HttpsError("aborted",
        "O bem foi alterado após a criação desta requisição. Revise a proposta com os dados atuais antes de responder.");
    // Resolve alvo existente ou prepara criação; helper só lê nesta etapa.
    const alvo = aprovar && req.novo_nome
      ? await prepararResumoAlvoTx(tx, idResumoAlvo, req.novo_nome) : null;
    const alteracoes = calcularAlteracoesPermitidas(req, bem);
    if (alvo) {
      if (alvo.novo) tx.create(alvo.ref, alvo.dados);
      alteracoes.id_resumo_bem_patrimonial = alvo.ref.id;
      alteracoes.nome_equipamento = alvo.dados.nome;
    }
    // M-02: quando novo_id_local está presente, propagar predio/andar/sala denormalizados
    if (aprovar && req.novo_id_local) {
      const localSnap = await tx.get(
        admin.firestore().collection("Local").doc(req.novo_id_local)
      );
      if (!localSnap.exists)
        throw new HttpsError("not-found", "Local referenciado na requisição não encontrado.");
      const local = localSnap.data()!;
      // Atualizar as projeções denormalizadas conforme Seção 5 (Espelhamento de Local)
      alteracoes.id_local = req.novo_id_local;
      alteracoes.predio = local.predio;
      alteracoes.andar = local.andar;
      alteracoes.sala = local.sala;
    }
    if (aprovar) {
      tx.update(bemRef, { ...alteracoes, versao: bem.versao + 1 });
      registrarHistoricoBemTx(tx, bemRef, bem, alteracoes, request.auth!.uid);
    }
    tx.update(reqRef, {
      status: aprovar ? "aprovada" : "rejeitada",
      novo_id_resumo_bem_patrimonial: alvo?.ref.id ?? null,
      respondida_em: admin.firestore.FieldValue.serverTimestamp(),
      id_usuario_respondente: request.auth!.uid, justificativa_resposta: justificativa,
    });
    tx.delete(admin.firestore().collection("Locks_Requisicao_Patrimonio").doc(`bem_edicao_${bemRef.id}`));
    return { status: aprovar ? "aprovada" : "rejeitada" };
  });
});

\end{lstlisting}

\begin{lstlisting}[language=TypeScript, title={criarRequisicaoAdicaoBem com lock determinístico (C2)}]
export const criarRequisicaoAdicaoBem = onCall(async (request) => {
  await validarPermissao(request, ["Professor"]);

  const dados = request.data as {
    numeroPatrimonioProposto: string;
    estadoConservacaoProposto: string;
    idLocal: string;
    nomeResponsavelProposto: string;
    photoUrlProposta: string; // objeto verificado pelo backend
    idResumoBemPatrimonial?: string;
    nomeResumoProposto?: string;
    descricaoResumoProposta?: string;
    motivo: string;
  };

  await validarFotoProposta(dados.photoUrlProposta, request.auth!.uid);

  // C2: Lock determinístico por número de patrimônio proposto
  const lockId = `bem_adicao_${dados.numeroPatrimonioProposto}`;
  const lockRef = admin.firestore().collection("Locks_Requisicao_Patrimonio").doc(lockId);
  const reqRef  = admin.firestore().collection("Requisicao_Adicao_Bem_Patrimonial").doc();

  return admin.firestore().runTransaction(async (tx) => {
    const lockSnap = await tx.get(lockRef);
    if (lockSnap.exists) {
      throw new HttpsError("failed-precondition", "Já existe requisição pendente para este número de patrimônio.");
    }
    tx.set(lockRef, { criado_em: admin.firestore.FieldValue.serverTimestamp() });
    tx.set(reqRef, {
      numero_patrimonio_proposto: dados.numeroPatrimonioProposto,
      estado_conservacao_proposto: dados.estadoConservacaoProposto,
      photo_url_proposta: dados.photoUrlProposta,
      nome_responsavel_proposto: dados.nomeResponsavelProposto,
      id_local: dados.idLocal,
      id_resumo_bem_patrimonial: dados.idResumoBemPatrimonial ?? null,
      nome_resumo_proposto: dados.nomeResumoProposto ?? null,
      descricao_resumo_proposta: dados.descricaoResumoProposta ?? null,
      motivo: dados.motivo,
      status: "pendente",
      feita_em: admin.firestore.FieldValue.serverTimestamp(),
      id_usuario_solicitante: request.auth!.uid,
      respondida_em: null,
      id_usuario_respondente: null,
    });
    return { idRequisicao: reqRef.id };
  });
});

export const responderRequisicaoAdicaoBem = onCall(async (request) => {
  await validarPermissao(request, ["Chefe_Geral", "Gestor_Bens_Patrimoniais"], true);
  const { idRequisicao, aprovar, justificativa } = request.data as { idRequisicao: string; aprovar: boolean; justificativa: string };
  const reqRef = admin.firestore().collection("Requisicao_Adicao_Bem_Patrimonial").doc(idRequisicao);

  return admin.firestore().runTransaction(async (tx) => {
    const reqSnap = await tx.get(reqRef);
    if (!reqSnap.exists) throw new HttpsError("not-found", "Requisição não encontrada.");
    const req = reqSnap.data()!;
    if (req.status !== "pendente") throw new HttpsError("failed-precondition", "Requisição já respondida.");

    if (aprovar && !req.photo_url_proposta)
      throw new HttpsError("failed-precondition", "Foto obrigatória: sanear pedido legado.");

    // C2: Remove o lock ao responder (libera possível nova requisição futura)
    const lockRef = admin.firestore().collection("Locks_Requisicao_Patrimonio")
      .doc(`bem_adicao_${req.numero_patrimonio_proposto}`);
    tx.delete(lockRef);

    let idBemCriado: string | null = null;

    if (aprovar) {
      let idResumo = req.id_resumo_bem_patrimonial;

      // Se não havia resumo existente, cria atômica e transacionalmente o Resumo_Bem_Patrimonial
      if (!idResumo && req.nome_resumo_proposto) {
        const novoResumoRef = admin.firestore().collection("Resumo_Bem_Patrimonial").doc();
        tx.set(novoResumoRef, {
          nome: req.nome_resumo_proposto,
          descricao: req.descricao_resumo_proposta ?? "",
          criado_em: admin.firestore.FieldValue.serverTimestamp(),
        });
        idResumo = novoResumoRef.id;
      }

      if (!idResumo) throw new HttpsError("failed-precondition", "A aprovação exige vincular um resumo.");

      const bemRef = admin.firestore().collection("Bem_Patrimonial").doc();
      idBemCriado = bemRef.id;

      tx.set(bemRef, {
        id_resumo_bem_patrimonial: idResumo,
        nome_equipamento: req.nome_resumo_proposto ?? "Equipamento",
        numero_patrimonio: req.numero_patrimonio_proposto,
        estado_conservacao: req.estado_conservacao_proposto,
        id_local: req.id_local,
        photo_url: req.photo_url_proposta, // ausência rejeitada antes das escritas
        documento_dado_baixa_pdf_url: null,
        nome_responsavel_sei: req.nome_responsavel_proposto,
        status: "Ativo",
        cadastrado_em: admin.firestore.FieldValue.serverTimestamp(),
      });

      tx.update(reqRef, {
        status: "aprovada",
        respondida_em: admin.firestore.FieldValue.serverTimestamp(),
        id_usuario_respondente: request.auth!.uid,
        id_bem_patrimonial_se_aprovado: idBemCriado,
        justificativa_resposta: justificativa,
      });
    } else {
      tx.update(reqRef, {
        status: "rejeitada",
        respondida_em: admin.firestore.FieldValue.serverTimestamp(),
        id_usuario_respondente: request.auth!.uid,
        justificativa_resposta: justificativa,
      });
    }

    return { status: aprovar ? "aprovada" : "rejeitada", idBemCriado };
  });
});

export const onResumoBemPatrimonialNomeAtualizado = onDocumentUpdated(
  "Resumo_Bem_Patrimonial/{resumoId}",
  async (event) => {
    const antes = event.data?.before.data();
    const depois = event.data?.after.data();
    if (!antes || !depois || antes.nome === depois.nome) return;

    const bens = await admin.firestore().collection("Bem_Patrimonial")
      .where("id_resumo_bem_patrimonial", "==", event.params.resumoId)
      .get();

    const batch = admin.firestore().batch();
    bens.docs.forEach((bem) => batch.update(bem.ref, {
      nome_equipamento: depois.nome,
    }));
    await batch.commit();
  }
);

export const onLocalAtualizado = onDocumentUpdated(
  "Local/{localId}",
  async (event) => {
    const antes = event.data?.before.data();
    const depois = event.data?.after.data();
    if (!antes || !depois) return;
    if (antes.predio === depois.predio && antes.andar === depois.andar && antes.sala === depois.sala) return;

    const bens = await admin.firestore().collection("Bem_Patrimonial")
      .where("id_local", "==", event.params.localId)
      .get();

    const batch = admin.firestore().batch();
    bens.docs.forEach((bem) => batch.update(bem.ref, {
      predio: depois.predio,
      andar: depois.andar,
      sala: depois.sala,
    }));
    await batch.commit();
  }
);

export const ingressarEmTurmaPorCodigo = onCall(async (request) => {
  await validarPermissao(request, ["Aluno"]);
  const { codigoTurma } = request.data as { codigoTurma: string };

  const turmaSnap = await admin.firestore().collection("Turma")
    .where("codigo_turma", "==", codigoTurma)
    .where("status", "==", "Ativo")
    .limit(1).get();
  if (turmaSnap.empty) throw new HttpsError("not-found", "Turma não encontrada ou arquivada.");
  const turmaDoc = turmaSnap.docs[0];
  const turma = turmaDoc.data();

  return admin.firestore().runTransaction(async (tx) => {
    // Checar capacidade
    const alunosSnap = await tx.get(
      turmaDoc.ref.collection("Alunos")
    );
    if (alunosSnap.size >= turma.capacidade) {
      throw new HttpsError("failed-precondition", "Turma lotada.");
    }

    // Checar se já está matriculado
    const jaMatriculado = alunosSnap.docs.some(
      (d) => d.id === request.auth!.uid
    );
    if (jaMatriculado) {
      throw new HttpsError("already-exists", "Você já está nesta turma.");
    }

    // P14: Checar se foi removido anteriormente
    const historico = await tx.get(
      turmaDoc.ref.collection("HistoricoAlunos")
        .where("id_aluno", "==", request.auth!.uid)
        .where("tipo", "==", "exclusao_aluno")
        .limit(1)
    );
    if (!historico.empty) {
      throw new HttpsError("failed-precondition",
        "Você foi removido desta turma pelo professor. O re-ingresso requer convite explícito.");
    }

    // Matricular
    tx.set(turmaDoc.ref.collection("Alunos").doc(request.auth!.uid), {
      id_aluno: request.auth!.uid,
      ingressou_em: admin.firestore.FieldValue.serverTimestamp(),
    });
    tx.set(turmaDoc.ref.collection("HistoricoAlunos").doc(), {
      id_aluno: request.auth!.uid,
      tipo: "inclusao_aluno",
      timestamp: admin.firestore.FieldValue.serverTimestamp(),
    });

    return { idTurma: turmaDoc.id, nomeTurma: turma.nome_turma };
  });
});
\end{lstlisting}

